\documentclass[10pt,a4paper]{amsart}
\usepackage[margin=19mm]{geometry}
\usepackage{amsmath,amssymb,mathtools}
\usepackage[scr=rsfs,cal=euler]{mathalfa}
\usepackage{xfrac}
\usepackage[unicode,hidelinks,bookmarks=false]{hyperref}
\newcommand{\ie}{i.e.\ }
\newcommand{\cf}{cf.\ }
\newcommand{\resp}{respectively\ }
\newcommand{\Subsection}{\subsection}
\providecommand{\nobreakdash}{\nobreak\hbox{-}}
\setlength{\emergencystretch}{2em}
\multlinegap0pt
\usepackage{tikz}
\usepackage{xcolor}
\usepackage{booktabs}
\usepackage{multirow}
\usepackage{amssymb}
\usepackage{mathtools}
\usepackage{algorithm}\usepackage{algpseudocode}
\usepackage{bm}
\usepackage{tcolorbox}
\newtheorem{restatable}{Restatable}
\newcommand{\bPhi}{\mathbf{\Phi}}
\newcommand{\bPsi}{\mathbf{\Psi}}
\title{Efficient Personalization of Generative Models via Optimal Experimental Design}
\author{Guy Schacht, Ziyad Sheebaelhamd, Riccardo De Santi, Mojm\'ir Mutn\'y, Andreas Krause}
\date{Source: arXiv:2512.19057v2 (CC BY 4.0)}
\begin{document}
\maketitle

\noindent\emph{Source and licence.} Efficient Personalization of Generative Models via Optimal Experimental Design. Version arXiv:2512.19057v2 (CC BY 4.0), distributed by arXiv under the Creative Commons Attribution 4.0 International licence, http://creativecommons.org/licenses/by/4.0/, as recorded in the arXiv licence metadata for this exact version and shown on the abstract page https://arxiv.org/abs/2512.19057v2. Full licence notice: \texttt{licenses/arxiv-2512.19057-CC-BY-4.0.txt}.\par\smallskip

\noindent\textit{Editorial scope.} Counted unit: the article's theorem FIMTruncatedVsStandard (source0.tex lines 1972--1989). The article's complete original proof of it is delivered with it (source0.tex lines 1991--2091, one \texttt{proof} environment of 101 lines). No mathematics is written, repaired or re-derived: the statement and the proof are the article's own lines, in the article's own notation, and no retained line is substituted or re-flowed.  No further source-local context is needed: every cross-reference of every retained line resolves inside the two delivered spans.  Nothing was omitted from any retained span, so the package carries no omission record.  No retained line contains a citation, so the wrapper carries no bibliography and no bibliography entry.  No retained line contains a figure environment, an image-inclusion command or drawing code, so no image is delivered and the package declares no asset dependency.  Editorial formatting only, in addition: an AMS \texttt{amsart} preamble that restates the article's own declarations and macros used by the retained lines, namely the environments \texttt{restatable} on separate counters, together with the standard packages those lines need.  The retained environments print in this document as Restatable 1 in order of appearance, whereas the article prints the same items with the numbering of its own full document.  The complete native source of the article as distributed in the arXiv e-print for this exact version is bundled unchanged as \texttt{sources/191443/source0.tex}.
\begin{restatable}[Comparison of Approximated FIMs]{theorem}{FIMTruncatedVsStandard}
\label{thm:fim_comparison}
Let $\mathcal{T}^K = \{ (\tau^1_t, \dots, \tau^K_t) \}_{t=1}^T$ be a set of $T \times K$ trajectories.
For the standard (state-based) feedback model, let $\phi(s^q_{t,h})$ be the feature vector for the state $s^q_{t,h}$. The approximated Fisher Information Matrix is
\begin{align*}
\tilde{I}^{\text{state}}(\mathcal{T}^K)
&= \sum_{t=1}^T \sum_{h=1}^H \Bigg( \frac{1}{K} \sum_{q=1}^K \phi(s^q_{t,h})\phi(s^q_{t,h})^\top \\
&\quad - \frac{1}{K^2} \sum_{q,q'=1}^K \phi(s^q_{t,h})\phi(s^{q'}_{t,h})^\top \Bigg).
\end{align*}
For the truncated trajectory feedback model, assume the perfect decomposition condition holds, such that the feature representation of the $q$-th trajectory in episode $t$ truncated at timestep $h$ is $\psi^q_{t,h} = \sum_{h'=1}^h \phi(s^q_{t,h'})$. The corresponding approximated Fisher Information Matrix is
\begin{align*}
\tilde{I}^{\text{trunc}}(\mathcal{T}^K)
&= \sum_{t=1}^T \sum_{h=1}^H \Bigg( \frac{1}{K} \sum_{q=1}^K \psi^q_{t,h}(\psi^q_{t,h})^\top \\
&\quad - \frac{1}{K^2} \sum_{q,q'=1}^K \psi^q_{t,h}(\psi^{q'}_{t,h})^\top \Bigg).
\end{align*}
Then, under the perfect decomposition condition,
\[ \tilde{I}^{\text{trunc}}(\mathcal{T}^K) \succeq \frac{1}{4} \cdot \tilde{I}^{\text{state}}(\mathcal{T}^K). \]
\end{restatable}

\begin{proof}[Proof of Theorem \ref{thm:fim_comparison}]
Let $\mathcal{T}^K = \{ (\tau^1_t, \dots, \tau^K_t) \}_{t=1}^T$ be the set of trajectories.
We define two approximated Fisher Information Matrices based on the uniform preference assumption ($p \approx 1/K$).

First, the FIM for the state-based feedback model, denoted $I^{\text{state}}_{\text{approx}}(\mathcal{T}^K)$, uses features $\phi(s^q_{t,h})$ from individual states:
    \begin{align*}
    I^{\text{state}}_{\text{approx}}(\mathcal{T}^K)
    &= \sum_{t=1}^T \sum_{h=1}^H \Bigg( \frac{1}{K} \sum_{q=1}^K \phi(s^q_{t,h})\phi(s^q_{t,h})^\top \\
    &\quad - \frac{1}{K^2} \sum_{q,q'=1}^K \phi(s^q_{t,h})\phi(s^{q'}_{t,h})^\top \Bigg).
    \end{align*}

Next, the FIM for the truncated trajectory feedback model, $I^{\text{trunc}}_{\text{approx}}(\mathcal{T}^K)$, under the perfect decomposition condition, uses the sum-decomposed features $\psi^q_{t,h} = \sum_{h'=1}^h \phi(s^q_{t,h'})$:
    \begin{align*}
    I^{\text{trunc}}_{\text{approx}}(\mathcal{T}^K)
    &= \sum_{t=1}^T \sum_{h=1}^H \Bigg( \frac{1}{K} \sum_{q=1}^K \psi^q_{t,h}(\psi^q_{t,h})^\top \\
    &\quad - \frac{1}{K^2} \sum_{q,q'=1}^K \psi^q_{t,h}(\psi^{q'}_{t,h})^\top \Bigg).
    \end{align*}

Let $\bPhi_{t,q} \in \mathbb{R}^{H \times d}$ be the matrix whose $h$-th row is $\phi(s^q_{t,h})^\top$. That is,
\[ \bPhi_{t,q} = \begin{pmatrix} \phi(s^q_{t,1})^\top \\ \phi(s^q_{t,2})^\top \\ \vdots \\ \phi(s^q_{t,H})^\top \end{pmatrix}. \]
The sum of outer products over the horizon $H$ for the state-based model is
\begin{align*}
\sum_{h=1}^H \phi(s^q_{t,h})\phi(s^q_{t,h})^\top &= \bPhi_{t,q}^\top I_H \bPhi_{t,q},
\end{align*}
where $I_H$ is the $H \times H$ identity matrix. Similarly, the sum of cross-products is
\begin{align*}
\sum_{h=1}^H \phi(s^q_{t,h})\phi(s^{q'}_{t,h})^\top &= \bPhi_{t,q}^\top I_H \bPhi_{t,q'}.
\end{align*}
Thus, $\tilde{I}^{\text{state}}(\mathcal{T}^K)$ can be rewritten as:
\begin{align*}
\tilde{I}^{\text{state}}(\mathcal{T}^K)
&= \sum_{t=1}^T \Bigg( \frac{1}{K} \sum_{q=1}^K \bPhi_{t,q}^\top I_H \bPhi_{t,q} \\
&\quad - \frac{1}{K^2} \sum_{q,q'=1}^K \bPhi_{t,q}^\top I_H \bPhi_{t,q'} \Bigg).
\end{align*}

For the truncated trajectory model (under perfect decomposition), let $\bPsi_{t,q} \in \mathbb{R}^{H \times d}$ be the matrix whose $h$-th row is $(\psi^q_{t,h})^\top = \left(\sum_{h'=1}^h \phi(s^q_{t,h'})\right)^\top$.
Let $S \in \mathbb{R}^{H \times H}$ be the lower triangular matrix of ones, i.e., $S_{ij}=1$ if $i \ge j$ and $S_{ij}=0$ if $i < j$. For example, if $H=3$:
\[ S = \begin{pmatrix} 1 & 0 & 0 \\ 1 & 1 & 0 \\ 1 & 1 & 1 \end{pmatrix}. \]
The cumulative sum structure means $\bPsi_{t,q} = S \bPhi_{t,q}$.
The sum of outer products over the horizon $H$ for the truncated model is
\begin{align*}
\sum_{h=1}^H \psi^q_{t,h}(\psi^q_{t,h})^\top
&= \bPsi_{t,q}^\top \bPsi_{t,q}
 = (S \bPhi_{t,q})^\top (S \bPhi_{t,q}) \\
&= \bPhi_{t,q}^\top S^\top S \bPhi_{t,q}.
\end{align*}
Similarly, the sum of cross-products is
\begin{align*}
\sum_{h=1}^H \psi^q_{t,h}(\psi^{q'}_{t,h})^\top
&= \bPsi_{t,q}^\top \bPsi_{t,q'} \\
&= \bPhi_{t,q}^\top S^\top S \bPhi_{t,q'}.
\end{align*}
Let $M = S^\top S$. This is an $H \times H$ symmetric positive definite matrix. For $H=3$,
\begin{align*}
M &= S^\top S \\
&= \begin{pmatrix} 1 & 1 & 1 \\ 0 & 1 & 1 \\ 0 & 0 & 1 \end{pmatrix}
   \begin{pmatrix} 1 & 0 & 0 \\ 1 & 1 & 0 \\ 1 & 1 & 1 \end{pmatrix} \\
&= \begin{pmatrix} 3 & 2 & 1 \\ 2 & 2 & 1 \\ 1 & 1 & 1 \end{pmatrix}.
\end{align*}
Thus, $\tilde{I}^{\text{trunc}}(\mathcal{T}^K)$ can be expressed in terms of $\bPhi_{t,q}$ and $M$:
\begin{align*}
\tilde{I}^{\text{trunc}}(\mathcal{T}^K)
&= \\
&\quad \sum_{t=1}^T \Bigg(
\frac{1}{K} \sum_{q=1}^K \bPhi_{t,q}^\top M \bPhi_{t,q} \\
&\qquad - \frac{1}{K^2} \sum_{q,q'=1}^K \bPhi_{t,q}^\top M \bPhi_{t,q'} \Bigg).
\end{align*}
Let
\[
X_t = \begin{pmatrix} \bPhi_{t,1}^\top & \dots & \bPhi_{t,K}^\top \end{pmatrix}^\top \in \mathbb{R}^{(KH) \times d}.
\]
Let $J_K = \frac{1}{K} \mathbf{1}_K \mathbf{1}_K^\top$ be the $K \times K$ matrix of all $1/K$.
Let $I_K$ be the $K \times K$ identity matrix.
The FIM expressions can be written compactly using Kronecker products $\otimes$:
\begin{align*}
\tilde{I}^{\text{state}}(\mathcal{T}^K) &= \sum_{t=1}^T X_t^\top \left( \frac{1}{K} (I_K - J_K) \otimes I_H \right) X_t \\
\tilde{I}^{\text{trunc}}(\mathcal{T}^K) &= \sum_{t=1}^T X_t^\top \left( \frac{1}{K} (I_K - J_K) \otimes M \right) X_t
\end{align*}
Since $M = S^\top S$ is positive definite (as $S$ is invertible), its eigenvalues are positive. Let $\lambda_{min}(M)$ be the smallest eigenvalue of $M$. Then $M \succeq \lambda_{min}(M) I_H$.
The matrix $I_K - J_K$ is positive semidefinite (it's proportional to a projection matrix).
Therefore, using properties of Kronecker products and Loewner order:
\begin{align*}
(I_K - J_K) \otimes M
&\succeq (I_K - J_K) \otimes (\lambda_{min}(M) I_H) \\
&= \lambda_{min}(M) (I_K - J_K) \otimes I_H
\end{align*}
Multiplying by $1/K$ and summing over $t$ after pre- and post-multiplying by $X_t^\top$ and $X_t$:
\begin{align*}
&\sum_{t=1}^T X_t^\top \left( \frac{1}{K} (I_K - J_K) \otimes M \right) X_t \\
&\quad \succeq \lambda_{min}(M) \sum_{t=1}^T X_t^\top \left( \frac{1}{K} (I_K - J_K) \otimes I_H \right) X_t
\end{align*}
This shows $\tilde{I}^{\text{trunc}}(\mathcal{T}^K) \succeq \lambda_{min}(M) \cdot \tilde{I}^{\text{state}}(\mathcal{T}^K)$.
The eigenvalues of $M=S^\top S$ are known to be
\[
\lambda_k(M) = \frac{1}{4 \sin^2\left(\frac{(2k-1)\pi}{2(2H+1)}\right)}, \quad k=1, \dots, H.
\]
Since $\sin^2(x) \le 1$ for any $x$, the minimum eigenvalue $\lambda_{min}(M)$ is lower-bounded by $1/4$.
Therefore, we can state the result using this constant lower bound:
\[ \tilde{I}^{\text{trunc}}(\mathcal{T}^K) \succeq \frac{1}{4} \cdot \tilde{I}^{\text{state}}(\mathcal{T}^K) \]
This completes the proof.
\end{proof}

\end{document}
